\section{Discussion}
\label{sec:discussion}
\subsection{On the Baseline Comparison}
\label{sec:baseline_discussion}
Table~\ref{tab:baseline} (Section~\ref{sec:baseline}) situates the proposed method relative to three existing Automotive Ethernet IDS approaches spanning the full spectrum of label requirements, from fully supervised~\cite{jeong2021convolutional} to semi-supervised~\cite{shibly2023feature} to unsupervised~\cite{jeong2023aero}. Among these, AERO~\cite{jeong2023aero} is the most directly comparable in terms of label efficiency, as it is likewise trained without attack labels; the meaningful distinction is therefore not unsupervised learning per se, but \emph{what} the unsupervised model is trained on. AERO's feature extractor constructs three engineered, protocol-informed features prior to learning, whereas the proposed method operates directly on raw packet bytes. This distinction matters in a multi-domain, multi-protocol setting such as CarDS: a feature extractor tailored to one protocol's semantics (e.g., AVTP or SOME/IP) may require redesign when applied to a different protocol or network segment, whereas the proposed raw-byte representation generalizes across segments by construction. Quantitatively, this label-efficiency advantage manifests as a per-category complementarity in the baselines' failure modes: AERO, the only other label-free method, drops to near-random on Emergency/VLAN\,4 Nmap scans (AUROC 0.50) despite handling RTP replacement reasonably (0.89), whereas the semi-supervised classifier inverts this profile, remaining near-perfect on scans (AUROC $\approx$1.0) but weakest on RTP replacement (0.74--0.84). The proposed method attains strong detection on both categories (RTP 0.978, Nmap 0.9998) without any attack labels or protocol-specific features, which we attribute to a single raw-byte representation that generalizes across heterogeneous attack types where hand-crafted or partially labeled features do not.

\subsection{Analysis of RTP Stream Replacement Detection}
Of the two attack categories evaluated, RTP stream replacement presents a structurally distinct detection challenge because it preserves packet-level structure --- valid RTP headers, consistent inter-arrival timing, and comparable payload lengths --- while substituting only the video payload content. This narrows the reconstruction error margin between normal and anomalous traffic, which is reflected in the threshold-independent separability reported in Section~\ref{sec:main_results} (AUPRC 0.9076 vs.\ 0.9998 for Nmap-based attacks). The per-file variation in F1-score within the category (0.8265--0.9217) further suggests that detectability is partly governed by the statistical distance between the injected video content and the training distribution of normal camera traffic, independent of the attack mechanism itself. Attack files whose replacement content is statistically closer to the normal distribution present a harder detection target --- an inherent property of reconstruction-based anomaly detection that implies evasion difficulty varies with the attacker's choice of replacement payload.

A notable observation is that time-aware positional encoding (Section~\ref{sec:pe_ablation}) improves RTP detection by AUPRC $+0.0401$ despite RTP stream replacement deliberately \emph{preserving} inter-packet timing, implying the gain does not arise from direct detection of timing anomalies. We offer two non-exclusive hypotheses. First, conditioning the Transformer's self-attention on real IPI data may yield a more precise learned representation of normal traffic globally --- including aspects unrelated to timing irregularities --- such that the content stream becomes more discriminative for RTP payload reconstruction as a secondary benefit. Second, despite the attacker's intent to preserve timing, the frame-replacement mechanism may itself introduce subtle inter-arrival perturbations at the moment of stream substitution, which a model conditioned on raw timing data may be more sensitive to than a rank-only sinusoidal baseline. Distinguishing between these hypotheses would require attention-weight analysis or a controlled ablation that perturbs timing while holding payload content fixed; we leave this to future work.

\subsection{Practical Deployment Considerations}
The proposed method is designed with practical automotive deployment in mind. With time-aware positional encoding, the model comprises approximately 897.6K trainable parameters (excluding the protocol-aware RTP header stream evaluated in Section~\ref{sec:ablation}), an increase of fewer than 17K parameters ($<$2\%) over the sinusoidal-PE backbone (880.8K), since the IPI projection consists of only two small linear layers. This corresponds to a memory footprint of approximately 3.4\,MB in 32-bit precision, remaining well within the constraints of automotive-grade embedded hardware. Training is performed exclusively on normal traffic captured during vehicle operation, requiring no attack simulation or manual annotation.

A practical consideration in sliding-window inference is the latency introduced by the requirement to buffer $N = 64$ consecutive packets before issuing a detection decision. In the CarDS dataset, the average inter-packet interval on the evaluated network segments is on the order of microseconds to milliseconds, resulting in a window accumulation time of well under one second under normal traffic conditions. This latency is acceptable for most automotive intrusion detection use cases, in which near-real-time rather than hard-real-time detection is sufficient. Nevertheless, reducing the window size while maintaining detection performance remains an open direction for latency-sensitive deployments, and is left for future investigation.

\subsection{Limitations}
\label{sec:limitations}

The evaluation is drawn entirely from a single 2020 production EV, and although the attack set spans 23 configurations, these reduce to only two underlying mechanisms -- Nmap-based scanning/spoofing and RTP stream replacement -- leaving other attack classes (e.g., denial-of-service, firmware/OTA tampering, and the excluded REST API trace) untested. The 0.01\% exclusion threshold also means behavior against deliberately low-footprint, stealthy attackers -- arguably the most realistic adversary -- remains unverified. Broader validation across vehicles, topologies, and attack families is needed.

Sliding windows use stride 1, so consecutive windows share 63 of 64 packets. Precision, recall, and F1 treat each window as an independent sample, which is standard in window-based intrusion detection but does not hold strictly under this overlap; reported metrics are best read as characterizing detection over the continuous stream rather than over independent trials.

Threshold selection is the one point in the pipeline where labeled data is used. Each percentile candidate ($P_{95}$--$P_{99.9}$) is computed purely from the normal validation distribution; only the choice among candidates uses labeled attacks (Section~\ref{sec:scoring}), drawn from the same 27 files used for final reporting rather than a disjoint split. This dependency has limited practical consequence, for two related reasons. Table~\ref{tab:threshold} shows that overall and RTP F1 remain broadly stable from $P_{95}$ through $P_{99.5}$, and the threshold-independent AUROC (0.9974) and AUPRC (0.9972) confirm this reflects genuine separability rather than a favorable threshold pick, so nearly any reasonable choice in the swept range would have performed comparably; only the extreme $P_{99.9}$ collapses. This same stability is also what makes the threshold robust to deployment-time drift: validation-set thresholds are not guaranteed to transfer exactly to deployment traffic, since the reconstruction error distribution may shift due to vehicle aging, software updates, or driving conditions absent from training data, yet the sinusoidal-PE baseline's RTP detection collapses almost entirely with a single percentile shift beyond its optimum, whereas the time-aware model's stable region extends substantially further, giving a wider safety margin against such drift. Even so, a fully label-free calibration procedure, or a held-out calibration split disjoint from the test files, would close the labeling gap more rigorously.

A more fundamental limitation is that the method assumes packet content carries learnable structure. Under link-layer or transport-layer encryption (e.g., MACsec, IPsec), byte content becomes effectively high-entropy and reconstruction-based detection loses its signal, regardless of traffic class. Pervasive per-packet encryption remains uncommon in current in-vehicle networks, however, mirroring CAN's own long-standing lack of encryption under tight compute and real-time budgets. This makes the assumption reasonable for the segments evaluated here, but it is a boundary condition rather than a guarantee, and should be revisited as automotive encryption adoption evolves.